\documentclass[12pt,a4paper]{article}

\usepackage[spanish,es-noquoting]{babel}
\usepackage[utf8]{inputenc}
\usepackage[T1]{fontenc}
\usepackage{lmodern}
\usepackage[a4paper,margin=2.5cm]{geometry}
\usepackage{setspace}
\usepackage{parskip}
\usepackage{hyperref}
\usepackage{xcolor}
\usepackage{enumitem}
\usepackage{booktabs}
\usepackage{longtable}
\usepackage{array}
\usepackage{fancyhdr}
\usepackage{titlesec}
\usepackage{graphicx}

\hypersetup{
    colorlinks=true,
    linkcolor=blue!60!black,
    urlcolor=blue!60!black,
    pdftitle={Informe Final - Control Simple de Inventario},
    pdfauthor={Nicolas Martinez y Martin Sanhueza}
}

\setstretch{1.15}
\pagestyle{fancy}
\fancyhf{}
\lhead{Control Simple de Inventario}
\rhead{Informe Final}
\cfoot{\thepage}

\titleformat{\section}{\large\bfseries\color{blue!60!black}}{\thesection}{0.6em}{}
\titleformat{\subsection}{\normalsize\bfseries\color{blue!50!black}}{\thesubsection}{0.6em}{}

\begin{document}

\begin{titlepage}
    \centering
    \vspace*{2cm}

    {\LARGE \textbf{Control Simple de Inventario}\par}
    \vspace{0.6cm}
    {\Large Informe final del Trabajo 2\par}
    \vspace{1.5cm}

    {\large \textbf{Integrantes}\par}
    \vspace{0.3cm}
    {\large Nicolas Martinez\par}
    {\large Martin Sanhueza\par}

    \vspace{1.2cm}
    {\large \textbf{Asignatura}\par}
    \vspace{0.3cm}
    {\large Desarrollo de Aplicaciones Empresariales\par}

    \vspace{1.2cm}
    {\large \textbf{Fecha}\par}
    \vspace{0.3cm}
    {\large 18 de abril de 2026\par}

\end{titlepage}

\tableofcontents
\newpage

\section{Introducción}

El proyecto \textit{Control Simple de Inventario} es una aplicación web orientada a la gestión de productos y movimientos de stock. La solución permite registrar productos, controlar entradas y salidas, visualizar alertas de inventario, revisar un historial de movimientos y consultar reportes ejecutivos. 

La aplicación fue desarrollada como una solución funcional de tipo empresarial, con frontend en React y backend en Node.js + Express, utilizando SQLite como base de datos persistente. El objetivo principal fue construir una aplicación que realmente funcionara de extremo a extremo, con datos almacenados de forma real y no solo simulada en memoria.

\section{Estructura de la aplicación}

La aplicación está organizada por capas para separar responsabilidades y facilitar su mantenimiento.

\subsection{Frontend}

La interfaz de usuario se encuentra en la carpeta \texttt{src/} y está compuesta por:

\begin{itemize}[leftmargin=1.2cm]
    \item \textbf{Componentes}: Header, SearchBar, ProductForm, ProductTable, MovementForm, MovementTable, AlertsPanel y ReportsPanel.
    \item \textbf{Hooks}: \texttt{useProducts.js}, que coordina estado, carga de datos y acciones principales.
    \item \textbf{Servicios}: \texttt{productService.js} y \texttt{storageService.js}, que concentran la lógica de negocio y comunicación con datos.
    \item \textbf{Utilidades}: validaciones y formateadores para fechas, cantidades y monedas.
    \item \textbf{Estilos}: CSS globales y componentes visuales, además de configuración de Tailwind CSS.
\end{itemize}

\subsection{Backend}

El backend se encuentra en la carpeta \texttt{backend/} y contiene:

\begin{itemize}[leftmargin=1.2cm]
    \item \texttt{server.js}: servidor Express con endpoints REST.
    \item \texttt{inventario.db}: base de datos SQLite que almacena productos y movimientos.
    \item \texttt{package.json}: scripts y dependencias del backend.
\end{itemize}

\newpage 

\subsection{Estructura general}

\begin{longtable}{>{\raggedright\arraybackslash}p{4cm} >{\raggedright\arraybackslash}p{10cm}}
\toprule
\textbf{Elemento} & \textbf{Función} \\
\midrule
Frontend React & Presentación, interacción y visualización de datos \\
Hook \texttt{useProducts} & Coordinación de estado y operaciones de la app \\
Servicios & Centralización de reglas de negocio y consumo de API \\
Backend Express & API REST y validación de persistencia \\
SQLite & Almacenamiento real de productos y movimientos \\
\bottomrule
\end{longtable}

\section{Configuración}

La solución final fue configurada con las siguientes tecnologías y decisiones:

\begin{itemize}[leftmargin=1.2cm]
    \item \textbf{React 18 + Vite} para el frontend.
    \item \textbf{Node.js + Express} para el backend.
    \item \textbf{SQLite} como base de datos local persistente.
    \item \textbf{TypeScript en migración progresiva}, manteniendo compatibilidad con el código existente.
    \item \textbf{Tailwind CSS} activado para apoyar el diseño visual.
    \item \textbf{ESLint} para revisión básica de calidad de código.
\end{itemize}

\subsection{Cómo se ejecuta}

Para correr el proyecto completo se requieren dos terminales:

\begin{enumerate}[leftmargin=1.2cm]
    \item En la carpeta \texttt{backend/}: \texttt{npm start}
    \item En la raíz del proyecto: \texttt{npm run dev}
\end{enumerate}

El backend queda disponible en \texttt{http://localhost:4000} y el frontend en \texttt{http://localhost:5173}.

\subsection{Persistencia}

La persistencia está implementada en SQLite mediante dos tablas principales:

\begin{itemize}[leftmargin=1.2cm]
    \item \textbf{products}: almacena código, nombre, cantidad, precio y fechas.
    \item \textbf{movements}: almacena entradas y salidas con cantidad anterior, cantidad nueva, motivo y referencia.
\end{itemize}

\newpage

\section{Funcionamiento}

La aplicación permite administrar inventario con un flujo completo:

\begin{enumerate}[leftmargin=1.2cm]
    \item Crear productos con validación de datos.
    \item Listar y buscar productos registrados.
    \item Editar o eliminar productos.
    \item Registrar movimientos de entrada y salida.
    \item Actualizar automáticamente el stock.
    \item Mostrar alertas de stock bajo o crítico.
    \item Generar reportes ejecutivos con métricas útiles.
\end{enumerate}

El flujo de datos es el siguiente:

\begin{center}
Usuario $\rightarrow$ Componente React $\rightarrow$ Hook $\rightarrow$ Servicio $\rightarrow$ API REST $\rightarrow$ SQLite $\rightarrow$ Respuesta $\rightarrow$ Re-render
\end{center}

\subsection{Funcionalidades destacadas}

\begin{itemize}[leftmargin=1.2cm]
    \item Control de stock en tiempo real.
    \item Historial de movimientos con antes y después.
    \item Alertas visuales para productos críticos.
    \item Reporte ejecutivo con indicadores de negocio.
    \item Persistencia real de datos.
\end{itemize}

\section{Forma de uso}

Para usar la aplicación se recomienda el siguiente flujo:

\begin{enumerate}[leftmargin=1.2cm]
    \item Abrir la aplicación en el navegador.
    \item Crear al menos un producto desde el formulario.
    \item Registrar entradas o salidas de stock.
    \item Revisar la tabla de productos para ver el stock actualizado.
    \item Consultar el historial de movimientos.
    \item Revisar alertas y reportes ejecutivos.
\end{enumerate}

En términos prácticos, el usuario puede:

\begin{itemize}[leftmargin=1.2cm]
    \item Agregar productos nuevos.
    \item Corregir información de productos existentes.
    \item Verificar si un producto está bajo de stock.
    \item Registrar una venta o una reposición.
    \item Analizar la trazabilidad de cada movimiento.
\end{itemize}

\section{Reflexión sobre el proceso de desarrollo}

\subsection{Qué intentamos hacer}

Nuestro objetivo fue construir una aplicación de inventario que no solo se viera bien, sino que realmente resolviera un problema completo: registrar productos, controlar stock, mantener historial y entregar información útil para tomar decisiones.

\subsection{Qué no funcionó y cómo lo enfrentamos}

Al comienzo, una parte importante de la lógica estaba basada solo en almacenamiento local. Eso hacía que la aplicación no tuviera persistencia real ni una separación clara entre frontend y backend. Para resolverlo, se migró el flujo hacia una API real con Express y SQLite.

También aparecieron problemas de consistencia entre la interfaz y las respuestas del backend, así como ajustes necesarios en el manejo asíncrono. Eso se corrigió organizando mejor los servicios, usando promesas y validando los resultados antes de actualizar la interfaz.

\subsection{Qué logramos con apoyo de la IA}

Con apoyo de la IA logramos acelerar tareas como:

\begin{itemize}[leftmargin=1.2cm]
    \item Ordenar la estructura general de la aplicación.
    \item Redactar documentación técnica más clara.
    \item Detectar inconsistencias entre la interfaz y la API.
    \item Ajustar estilos visuales y mejorar legibilidad.
    \item Refinar el contenido para la entrega final y el inicio rápido.
\end{itemize}

\subsection{Qué logramos por nuestra cuenta}

Sin intervención directa de la IA logramos:

\begin{itemize}[leftmargin=1.2cm]
    \item Entender el funcionamiento general de la aplicación.
    \item Probar la app en el navegador y revisar su comportamiento real.
    \item Identificar necesidades de usabilidad y presentación.
    \item Tomar decisiones sobre qué debía quedar en la versión final y qué debía quedar como historial técnico.
\end{itemize}

\subsection{Qué parte del sistema no podríamos explicar con total claridad y por qué}

La parte que podría requerir más revisión para explicarla con total seguridad sería la configuración detallada de algunas herramientas de compilación y estilos, especialmente la combinación entre Tailwind CSS, PostCSS y la estructura de build del proyecto. También sería necesario revisar nuevamente ciertos detalles del servidor y de la persistencia para explicar cada paso de implementación con precisión absoluta.

Esto no significa que no entendamos el sistema, sino que existen detalles de configuración fina que conviene verificar antes de exponerlos en una evaluación oral.

\section{Qué aprendimos para la evaluación}

La evaluación del curso se centra en la experiencia de desarrollo y no en detalles finos de código. Por eso, lo más importante que aprendimos fue:

\begin{itemize}[leftmargin=1.2cm]
    \item Cómo construir una aplicación que realmente funcione de extremo a extremo.
    \item Por qué es importante separar presentación, lógica y persistencia.
    \item Qué errores son comunes al integrar frontend y backend.
    \item Cómo identificar problemas y corregirlos de forma progresiva.
    \item Que una aplicación útil no depende solo de que se vea bien, sino de que resuelva un flujo real.
\end{itemize}

\section{Conclusión}

El proyecto \textit{Control Simple de Inventario} cumple con el objetivo de la actividad porque integra interfaz, lógica, backend y persistencia real en SQLite. La aplicación permite registrar productos y movimientos, controlar el stock, generar alertas y revisar información útil para la gestión.

Además, el proceso de desarrollo dejó aprendizaje sobre construcción de aplicaciones funcionales, depuración, estructura del proyecto y documentación técnica. Con ello, la entrega final queda preparada tanto para demostración como para evaluación.

\section{Anexo breve: datos de entrega}

\begin{itemize}[leftmargin=1.2cm]
    \item Repositorio o ZIP de la aplicación: incluir junto a esta memoria.
    \item Video de demostración: máximo 5 minutos.
    \item Entrega escrita: este documento en PDF.
\end{itemize}

\end{document}
